\documentclass[10pt,a4paper]{amsart}
\usepackage[margin=19mm]{geometry}
\usepackage{amsmath,amssymb,mathtools}
\usepackage[scr=rsfs,cal=euler]{mathalfa}
\usepackage{xfrac}
\usepackage[unicode,hidelinks,bookmarks=false]{hyperref}
\newcommand{\ie}{i.e.\ }
\newcommand{\cf}{cf.\ }
\newcommand{\resp}{respectively\ }
\newcommand{\Subsection}{\subsection}
\providecommand{\nobreakdash}{\nobreak\hbox{-}}
\setlength{\emergencystretch}{2em}
\multlinegap0pt
\usepackage{amssymb}
\usepackage[shortlabels]{enumitem}
\newtheorem{theorem}{Theorem}
\title{An explicit family of 30 blocks meeting every 6-set of [60] in at least two points}
\author{Paulo Henrique Cunha Gomes}
\date{Source: arXiv:2601.04295v1 (CC BY 4.0)}
\begin{document}
\maketitle

\noindent\emph{Source and licence.} An explicit family of 30 blocks meeting every 6-set of [60] in at least two points. Version arXiv:2601.04295v1 (CC BY 4.0), distributed by arXiv under the Creative Commons Attribution 4.0 International licence, http://creativecommons.org/licenses/by/4.0/, as recorded in the arXiv licence metadata for this exact version and shown on the abstract page https://arxiv.org/abs/2601.04295v1. Full licence notice: \texttt{licenses/arxiv-2601.04295-CC-BY-4.0.txt}.\par\smallskip

\noindent\textit{Editorial scope.} Counted unit: the article's theorem thm:main (source0.tex lines 85--90). The article's complete original proof of it is delivered with it (source0.tex lines 92--117, one \texttt{proof} environment of 26 lines). No mathematics is written, repaired or re-derived: the statement and the proof are the article's own lines, in the article's own notation, and no retained line is substituted or re-flowed.  No further source-local context is needed: every cross-reference of every retained line resolves inside the two delivered spans.  Nothing was omitted from any retained span, so the package carries no omission record.  No retained line contains a citation, so the wrapper carries no bibliography and no bibliography entry.  No retained line contains a figure environment, an image-inclusion command or drawing code, so no image is delivered and the package declares no asset dependency.  Editorial formatting only, in addition: an AMS \texttt{amsart} preamble that restates the article's own declarations and macros used by the retained lines, namely the environments \texttt{theorem} on separate counters, together with the standard packages those lines need.  The retained environments print in this document as Theorem 1 in order of appearance, whereas the article prints the same items with the numbering of its own full document.  The complete native source of the article as distributed in the arXiv e-print for this exact version is bundled unchanged as \texttt{sources/192287/source0.tex}.
\begin{theorem}\label{thm:main}
Let $\mathcal{B}$ be the family of $30$ blocks constructed above. For every $S\subset[60]$ with
$|S|=6$, there exists $B\in\mathcal{B}$ such that $|S\cap B|\ge 2$.
Equivalently: every $6$-subset $S$ contains a pair $\{x,y\}\subset S$ that is contained in at least
one block $B\in\mathcal{B}$.
\end{theorem}

\begin{proof}
Fix $S\subset[60]$ with $|S|=6$.

\smallskip
\noindent\textbf{Case 1.} If $S$ contains two elements from the same base block $G_i$, then
$|S\cap G_i|\ge 2$ and we are done since $G_i\in\mathcal{B}$.

\smallskip
\noindent\textbf{Case 2.} Otherwise, $S$ meets each base block in at most one element. Since
$|S|=6$, the set $S$ meets exactly six distinct base blocks among $G_1,\dots,G_{10}$.

The ten base blocks are partitioned into five pairs
\[
(G_1,G_2),\ (G_3,G_4),\ (G_5,G_6),\ (G_7,G_8),\ (G_9,G_{10}).
\]
By the pigeonhole principle, among the six base blocks hit by $S$, at least two belong to the same
pair, say $(G_{2m-1},G_{2m})$. Thus $S$ contains one element
$x\in G_{2m-1}$ and one element $y\in G_{2m}$.

By construction, $x\in G_{2m-1}^{(u)}$ for some $u\in\{1,2\}$ and
$y\in G_{2m}^{(v)}$ for some $v\in\{1,2\}$. Therefore the recombined block
\[
B:=G_{2m-1}^{(u)}\cup G_{2m}^{(v)}
\]
belongs to $\mathcal{B}$ and contains both $x$ and $y$, hence $|S\cap B|\ge 2$.
\end{proof}

\end{document}
